%
%
%
%

\section{Introduction} \label{sec:intro}

%
%

%
%

Network data, that is, the encoding of connections between objects, is a natural form of information representation in many fields \cite{Porter2020}.
While the analysis of static networks or graphs is already a broad field of research in itself, there is often much information ignored. 
In particular, we are interested in the case of \textit{temporal networks} \cite{Holme2013,Holme2015}; that is, the case of a dynamical system represented by a network evolving over time. 
These networks can arise in many different cases, such as 
social networks~\cite{Skyrms2000}, 
disease spread dynamics~\cite{Husein2019}, 
manufacturer-supplier networks~\cite{Xu2019b}, 
power grid network~\cite{Schaefer2018}, and
transportation networks~\cite{DavidBoyce2012}. 
%
Many important characteristics of a dynamical network can be extracted from the data. These include source and rate of disease spread as well as predictions on future infections~\cite{Enright2018}, weak branches in supply chains and possible failures~\cite{Nuss2016, Xu2019b}, changes in infrastructure to avoid cascade failures in power grids~\cite{Soltan2014, Schaefer2018}, transportation network optimal routing (finding an optimal minimum time route between)~\cite{Bast2015}, fault analysis (detecting transportation disruptions) in transportation networks~\cite{Sugishita2020}, and flow pattern analysis (visualization)~\cite{Hackl2019}.


Temporal graph data is commonly represented using attributed information on the edges for the time intervals or instances in which the edges are active~\cite{Chen2016a,Huang2015}. 
Using this attributed information, we can represent the graph in several ways including edge labeling, snapshots, and static graph representations~\cite{Wang2019}. 
In this work we will first represent the data in the standard attributed (labeled) temporal graph structure and then use the graph snapshots approach, where our input is a sequence of static graphs $G_0, G_1, \ldots, G_n$.


The available tools for analyzing such data is often inspired by the tools from the static network community; for instance, 
centrality or flow measures~\cite{Borgatti2005}; temporal clustering for event detection~\cite{Crawford2018, You2021, Moriano2019}; and connectedness~\cite{Kempe2002}. 
However, these tools do not account for higher-dimensional structures (e.g., loops as a one-dimensional structure). 
It may be important to account for evolving higher-dimensional structures in temporal networks to understand the changing structure better. For example, a highly connected network may only have one connected component with no clear clusters, but the number of loops within the network may detect the change.

For this reason, we introduce a method for incorporating ideas from Topological Data Analysis (TDA) \cite{Dey2021,Munch2017} to encode more complex structure than can be seen in the standard graph tools. 
The mainstay of TDA is persistent homology, colloquially referred to as persistence, which  encodes structure by analyzing the changing shape of a simplicial complex (a higher dimensional generalization of a network) over a filtration (a nested sequence of subcomplexes $K_1 \subseteq K_2 \subseteq \cdots \subseteq K_n$).
This shape is measured via homology, a vector space encoding information about topological structure of the space. 
Different dimensions of homology measure different things: 0-dimensional homology encodes information about connected components; 1-dimensional homology encodes loops; 2-dimensional homology encodes voids. 
How these structures change over the filtration (e.g.~when a loop appears and then subsequently fills in) can be stored as a persistence diagram; that is, a collection of points in the upper half plane where a point at $(i,j)$ means that a structure appeared at $K_i$ but was lost at $K_j$. 


However, a major limitation of standard persistence for the temporal network input data is the requirement that inclusions only go one way. 
Over the evolution parameter, our temporal graphs might add or remove edges, and we would like to build a system that can account for this. 
Thus, to study the evolving higher dimensional structures within a temporal network, we will leverage zigzag persistence~\cite{Carlsson2010}, which allows for insertions and deletions of simplices at every step.
The same mathematical theorems that make it possible to represent the information of standard persistence in a persistence diagram or barcode can be used in the case of a zigzag filtration. 
Zigzag persistence tracks the formation and disappearance of homological structures through a persistence diagram as a two-dimensional summary diagram. 

The tools we employ will treat the graph input itself as a topological structure by using it to extract distance information between the vertices. 
For the work discussed here, we will assume the input networks are unweighted graphs and thus use the shortest path distance between vertices. 
However, this is by no means the only way to incorporate graph data for persistent homology; see \cite{Aktas2019} for a survey. 
Further, because we generally treat our networks as a metric space, our intuition is that the tools developed here are most applicable to data where vertices have a geometric component (such as networks where nodes come from geo-spatial data) rather than only combinatorial data (such as human proximity networks, even if this has to do with physical proximity), but we look forward to being proven wrong.

We thus apply our methods to two sample data sets. 
The first is transportation data from Great Britain \cite{Gallotti2015}, where we can use the zigzag persistent homology for understanding higher order structures in the system, and see behavior such as daily periodicity represented automatically in the diagram output. 
The second data comes from input data from time series analysis. 
We have previously used zigzag persistence to detect Hopf bifurcations  by taking as input a point cloud approximation of the underlying attractor to the time series \cite{Tymochko2020}. 
However, the computational limitations of zigzag persistence on large point clouds makes its use in this setting unwieldy. 
Thus, we combine these ideas with recently available network representations of time series \cite{McCullough2015} to 
%
instead encode the structure of the time series in a graph.
Previous work shows that measuring this structure using standard persistent homology can be used to differentiate between behaviors in the underlying dynamical system \cite{Myers2019,Myers2022}. 





%
%
%
%
%
%
%
%

%
%
%
%

%
%
%
%
%
%
%
%
%

%
%

%
%
%
%
%
%
%

%
%

%
%

%
%


%


\subsection{Organization}

We will start in Section~\ref{sec:background} with an introductory background on persistent homology and zigzag persistence.
Next, in Section~\ref{sec:method},  we overview the general pipeline for applying zigzag persistence to temporal graph data. We couple this explanation with a demonstrative toy example. 
In Section~\ref{sec:results} we  introduce the two systems we will study. 
The first is a dataset collected over a week of the Great Britain transportation system~\cite{Gallotti2015}. 
The second is an intermittent Lorenz system simulation, where we generate a temporal network through complex networks of sliding windows. %
Then we apply zigzag persistence to our two examples and show how the resulting persistence diagrams help visualize the underlying dynamics in comparison to standard temporal network analysis techniques.



